\documentclass[11pt,a4paper]{article}

%% ── Encoding & fonts ──────────────────────────────────────────────────────────
\usepackage[utf8]{inputenc}
\usepackage[T1]{fontenc}
\usepackage{lmodern}

%% ── Core packages ────────────────────────────────────────────────────────────
\usepackage{amsmath,amssymb,amsfonts,amsbsy}
\usepackage{amsthm}
\usepackage{mathrsfs}
\usepackage{bm}
\usepackage{geometry}
\usepackage{xcolor}
\usepackage{booktabs}
\usepackage{array}
\usepackage{longtable}
\usepackage{tabularx}
\usepackage{enumitem}
\usepackage{fancyhdr}
\usepackage{titlesec}
\usepackage{hyperref}
\usepackage{microtype}
\usepackage{float}
\usepackage{caption}
\usepackage{mdframed}
\usepackage{setspace}
\usepackage{graphicx}

%% ── Page geometry ────────────────────────────────────────────────────────────
\geometry{left=2.8cm,right=2.8cm,top=3.2cm,bottom=3.0cm,headheight=15pt}

%% ── Hyperref ─────────────────────────────────────────────────────────────────
\hypersetup{
  colorlinks=true, linkcolor=black, citecolor=black, urlcolor=black,
  pdftitle={Looking Up: A Projection-Residual Framework Without Compactified Extra Dimensions}
}

%% ── Section styling ──────────────────────────────────────────────────────────
\titleformat{\section}{\normalfont\large\bfseries}{\thesection.}{0.6em}{}[\vspace{-2pt}\rule{\textwidth}{0.4pt}]
\titleformat{\subsection}{\normalfont\normalsize\bfseries}{\thesubsection}{0.5em}{}
\titleformat{\subsubsection}{\normalfont\normalsize\itshape}{\thesubsubsection}{0.5em}{}

%% ── Header / footer ──────────────────────────────────────────────────────────
\pagestyle{fancy}
\fancyhf{}
\lhead{\footnotesize\textit{Looking Up}\textbar{}\ Fifth-Dimensional Momentum Framework}
\rhead{\footnotesize Rev.~1.0}
\lfoot{\footnotesize\today}
\rfoot{\footnotesize\thepage}
\renewcommand{\headrulewidth}{0.4pt}
\renewcommand{\footrulewidth}{0.3pt}

%% ── Notation macros ──────────────────────────────────────────────────────────
\newcommand{\pw}{p_w}                        % hidden-axis momentum scalar
\newcommand{\pvw}{\mathbf{p}_w}              % hidden-axis momentum vector context
\newcommand{\Lf}{\Lambda_5}                  % fifth-force coupling constant
\newcommand{\afive}{\mathbf{a}_5}            % fifth-force acceleration (vector)
\newcommand{\afiver}{a_5}                    % fifth-force acceleration (scalar/radial)
\newcommand{\Phifive}{\Phi_5}                % fifth-force potential
\newcommand{\Seff}{S_{\rm eff}}              % effective switching function
\newcommand{\Sfive}{S_5}                     % radial propagation function
\newcommand{\rcrit}{r_{\rm crit}}            % transition/switching radius
\newcommand{\vobs}{v_{\rm obs}}              % observed velocity
\newcommand{\vbar}{v_{\rm bar}}              % baryonic velocity
\newcommand{\vmod}{v_{\rm model}}            % model velocity
\newcommand{\aobs}{a_{\rm obs}}              % observed acceleration
\newcommand{\abar}{a_{\rm bar}}              % baryonic acceleration
\newcommand{\asnap}{a_5^{\rm snap}}          % snapshot fifth force
\newcommand{\atot}{a_5^{\rm total}}          % total fifth force
\newcommand{\rhat}{\hat{\mathbf{r}}}         % radial unit vector
\newcommand{\Bfield}{\mathcal{B}_5}          % background field
\newcommand{\vfive}{v_5}                     % fifth-force velocity contribution
\newcommand{\vN}{v_N}                        % Newtonian velocity
\newcommand{\InfoD}{\Delta\mathfrak{D}}      % old information delta
\newcommand{\UA}{U^A}                        % 5D four-velocity
\newcommand{\XA}{X^A}                        % 5D position
\newcommand{\PA}{P^A}                        % 5D momentum
\newcommand{\dS}{dS_5^2}                     % 5D interval
\newcommand{\Kw}{K_w}                        % hidden-axis kinetic energy

%% ── Key-result box environment ───────────────────────────────────────────────
\newmdenv[
  topline=true, bottomline=true, rightline=false, leftline=false,
  linewidth=1.2pt, linecolor=black,
  innertopmargin=8pt, innerbottommargin=8pt,
  innerleftmargin=4pt, innerrightmargin=4pt,
  skipabove=10pt, skipbelow=10pt
]{keybox}

\captionsetup{font=small,labelfont=bf}

%% ── Figure suggestion annotation ─────────────────────────────────────────────
\newcommand{\figsuggest}[1]{%
  \medskip
  \noindent\hspace{0pt}\textcolor{gray}{\small\itshape
    $\triangleright$\;Figure suggestion:\enspace #1}%
  \medskip}

%% ═════════════════════════════════════════════════════════════════════════════
\begin{document}
%% ═════════════════════════════════════════════════════════════════════════════

%% ── TITLE PAGE ───────────────────────────────────────────────────────────────
\thispagestyle{empty}
\vspace*{2.2cm}

\begin{center}
  \noindent\rule{\textwidth}{1.2pt}\vspace{1.0em}

  {\fontsize{24}{30}\selectfont\bfseries Looking Up}

  \vspace{0.6em}
  {\large\itshape
    A Fifth-Dimensional Momentum Framework\\[0.2em]
    for Galactic Rotation Without Dark Matter
  }

  \vspace{0.4em}
  {\normalsize
    A Projection-Residual Framework Without Compactified Extra Dimensions
  }

  \vspace{0.8em}
  \noindent\rule{\textwidth}{0.5pt}\vspace{0.8em}

  {\normalsize\textsc{Theoretical Division}}\\[0.3em]
  {\small\textit{Particle Physics and Cosmological Modelling Research Group}}

  \vspace{1.8em}
  {\small
  \begin{tabular}{rl}
    \textbf{Identifier:} & LOOKINGUP-1.0 \\[0.3em]
    \textbf{Revision:}   & 1.0 \\[0.3em]
    \textbf{Date:}       & \today \\[0.3em]
    \textbf{Status:}     & Theoretical framework draft \\
  \end{tabular}
  }

  \vspace{2.0em}
  \noindent\rule{\textwidth}{0.5pt}\vspace{0.8em}

  \begin{minipage}{0.86\textwidth}
  \small\textbf{Abstract.}\quad
  This paper develops an Inverse Kaluza-Klein Scale Symmetry framework as a
  projection-residual approach to galactic rotation, force behaviour, and
  scale-dependent observation.  Unlike standard Kaluza-Klein theory, the
  framework does not assume that an extra dimension is compactified into a
  microscopic circular coordinate.  Instead it reverses the inference
  direction: observed force-scale behaviour, hidden-active internal structure,
  and unresolved projection residuals are used to infer the existence of a
  missing scale permission.  The central claim is that ordinary spacetime
  observation $O_S = P_S(I)$ does not recover the full physical identity $I$,
  and that the residual $R_S = U_S(I) - E_S(O_S)$ encodes unresolved identity
  content at each observational scale.

  The $w$-coordinate introduced here is not a curled spatial dimension.
  It is a scale-depth degree of freedom associated with projection access,
  identity packing, and residual structure.  Within this framing,
  hidden-axis momentum $\pw = m\,dw/d\tau$ acts as a fifth-force charge whose
  radial projection follows a logarithmic coupling.  The resulting acceleration
  \[
    \afiver(r) = \Lf\ln\!\left(1+\frac{|\pw|}{p_0}\right)
    G(r)\frac{1}{r+r_c}
  \]
  naturally produces a nearly flat contribution to $v^2(r)$ at large galactic
  radii.  A sigmoid activation gate $G(r)$ suppresses the fifth force at all
  sub-galactic scales, protecting solar-system tests and molecular dynamics
  without patching.  Three interpretive schools, a complete numerical
  parameterisation protocol for the SPARC database, and precise falsification
  criteria are developed.  Dark-sector signatures are reinterpreted as
  gravitationally visible but electromagnetically unresolved identity content
  --- a projection gap rather than invisible matter.
  \end{minipage}

  \vspace{0.8em}
  \noindent\rule{\textwidth}{1.2pt}
\end{center}

%% ── Derivation summary ───────────────────────────────────────────────────────
\vspace{1.2em}
\begin{center}
\begin{minipage}{0.88\textwidth}
\small
\textbf{Step-by-step derivation of $\boldsymbol{\mathbf{a}_5}$.}\quad
The fifth-force acceleration follows from four sequential steps:

\begin{enumerate}[leftmargin=1.8em,itemsep=4pt]

\item \textbf{5D proper time.}\enspace
  The five-dimensional spacetime interval $\dS = -c^2dT^2 + d\mathbf{r}^2 + dw^2$
  implies a proper-time relation that includes motion along the hidden axis $w$:
  \[
    d\tau = dT\sqrt{1 - \frac{|\mathbf{v}|^2}{c^2} - \frac{v_w^2}{c_w^2}} .
  \]

\item \textbf{Hidden-axis momentum.}\enspace
  The projection of the 5D momentum vector onto the $w$-direction defines the
  fifth-force charge.  In the non-relativistic (galactic) limit:
  \[
    \pw = m\,\frac{dw}{d\tau} \;\approx\; m\,v_w .
  \]

\item \textbf{Logarithmic potential.}\enspace
  A logarithmic coupling $g(\pw) = \ln(1+|\pw|/p_0)$ is chosen for its
  saturation behaviour at large $|\pw|$ and vanishing at $|\pw|=0$.
  The resulting fifth-force potential is:
  \[
    \Phi_5(r) = -\Lf \ln\!\left(1+\frac{|\pw|}{p_0}\right)
    \ln\!\left(\frac{r_c}{r + r_c}\right) .
  \]

\item \textbf{Radial gradient + gate.}\enspace
  Differentiating with respect to $r$ and multiplying by the sigmoid activation
  gate $G(r) = 1/(1+e^{-k(r-r_0)})$:
  \[
    \boxed{\,
      \afive(r) = -\Lf\ln\!\left(1+\frac{|\pw|}{p_0}\right)
      \underbrace{\frac{1}{1+e^{-k(r-r_0)}}}_{G(r)}
      \frac{1}{r+r_c}\,\rhat
    \,}
  \]
  At large galactic radii, $G(r)\to 1$ and $r+r_c \approx r$, giving a
  $1/r$ force profile that produces flat rotation curves.
  At sub-galactic radii, $G(r)\to 0$ and the fifth force vanishes.

\end{enumerate}
\end{minipage}
\end{center}

\newpage
\tableofcontents
\newpage

%% ── Post-TOC derivation reference ────────────────────────────────────────────
\thispagestyle{fancy}
\vspace*{0.6cm}
\begin{center}
\begin{minipage}{0.90\textwidth}

\textbf{\large Derivation Reference: $\boldsymbol{\mathbf{a}_5}$ and School~I Curve Fitting}

\smallskip\rule{\textwidth}{0.5pt}\smallskip

\textbf{A.\ Full derivation chain for $\afive$.}

\smallskip
\noindent The fifth-force acceleration is derived in four steps from the
5D spacetime manifold, requiring no new force postulates --- only the
existence of a hidden spatial axis $w$ and the standard relation between
momentum and velocity along that axis.

\begin{enumerate}[leftmargin=1.8em,itemsep=6pt]

\item \textbf{5D interval and proper time.}
\[
  dS_5^2 = -c^2dT^2 + d\mathbf{r}^2 + dw^2
  \;\;\Rightarrow\;\;
  d\tau = dT\sqrt{1 - \tfrac{|\mathbf{v}|^2}{c^2} - \tfrac{v_w^2}{c_w^2}} .
\]

\item \textbf{Hidden-axis momentum as fifth-force charge.}
\[
  \pw \;=\; m\frac{dw}{d\tau} \;\approx\; mv_w
  \quad\text{(non-relativistic galactic limit).}
\]

\item \textbf{Logarithmic potential from coupling function $g(\pw) = \ln(1+|\pw|/p_0)$.}
\[
  \Phi_5(r) = -\Lf\ln\!\left(1+\tfrac{|\pw|}{p_0}\right)\ln\!\left(\tfrac{r_c}{r+r_c}\right).
\]

\item \textbf{Radial gradient and sigmoid gate $G(r) = 1/(1+e^{-k(r-r_0)})$.}
\[
  \boxed{
    \afive(r) = -\Lf\ln\!\left(1+\tfrac{|\pw|}{p_0}\right)
    \cdot G(r)
    \cdot \tfrac{1}{r+r_c}\,\rhat
  }
\]
\end{enumerate}

\medskip
\rule{\textwidth}{0.4pt}
\smallskip

\textbf{B.\ School~I (Kinematic Snapshot) curve-fitting protocol.}

\smallskip
\noindent
The kinematic snapshot model treats $|\pw(r)|$ as a radial proxy sourced
by the baryonic momentum distribution.  The complete curve-fitting derivation
for a single galaxy proceeds as follows.

\medskip
\noindent\textit{Step 1 — baryonic baseline.}\enspace
Take the baryonic contribution $\vbar^2(r) = GM_{\rm bar}(r)/r$ directly
from stellar mass models and H{\sc i} gas maps.  This is \emph{not} fitted;
it is an input.

\medskip
\noindent\textit{Step 2 — observed residual.}\enspace
Compute the residual acceleration that baryons alone cannot explain:
\[
  \Delta a(r) = \frac{\vobs^2(r) - \vbar^2(r)}{r} .
\]

\medskip
\noindent\textit{Step 3 — fifth-force proxy.}\enspace
Adopt a physically motivated proxy for hidden-axis momentum.
The three canonical choices, in order of increasing complexity, are:
\begin{align*}
  |\pw(r)| &\propto m_* v_{\rm circ}(r)
  & &\text{(circular velocity proxy),} \\
  |\pw(r)| &\propto \Sigma(r)\,v_{\rm circ}(r)
  & &\text{(surface-density proxy),} \\
  |\pw(r)| &\propto \rho(r)\,v(r)
  & &\text{(momentum-flux proxy).}
\end{align*}

\medskip
\noindent\textit{Step 4 — parameter optimisation.}\enspace
Minimise over the universal parameter vector
$\Theta = \{\Lf,\,p_0,\,r_c,\,r_0,\,k\}$:
\[
  \chi^2(\Theta) = \sum_{j=1}^{N}
  \left[\frac{\vobs(r_j) - \vmod(r_j;\,\Theta)}{\sigma_{v,j}}\right]^2,
  \quad
  \vmod^2(r) = \vbar^2(r)
  + \Lf\,r\ln\!\left(1+\tfrac{|\pw(r)|}{p_0}\right)\Sfive(r;\Theta) .
\]

\medskip
\noindent\textit{Step 5 — population consistency gate.}\enspace
Repeat for $N_{\rm gal} \geq 10$ galaxies.  Accept the model only if
$\{\Lf,\,p_0,\,r_c\}$ remain within a factor of $\sim$3 across the sample.
Per-galaxy quantities $\{r_0,\,k,\,|\pw(r)|\}$ may vary but must
correlate with observable kinematics.

\medskip
\noindent\textit{Step 6 — falsification check.}\enspace
If Step~5 fails --- i.e., the universal parameters show order-of-magnitude
variation between galaxies with no kinematic predictor --- the School~I
model is falsified and must be reclassified as a disguised halo fit.

\medskip
\rule{\textwidth}{0.5pt}
\end{minipage}
\end{center}
\newpage

%%═══════════════════════════════════════════════════════════════════════════════
\section{Introduction: The Rotation Curve Problem}\label{sec:intro}
%%═══════════════════════════════════════════════════════════════════════════════

In a purely Newtonian framework, stars and gas in the outer disk of a spiral
galaxy should orbit analogously to the planets of the outer solar system: as
enclosed mass decreases with radius, orbital velocities should fall.  They do
not.  From the Andromeda Galaxy to low-surface-brightness dwarfs, the observed
outer-disk velocity profile is approximately flat where theory predicts it
should decline.  The discrepancy is systematic, persistent, and large --- at
outer-disk radii, visible-matter predictions can fall short of observed
velocities by factors of two or more.

The dominant resolution invokes dark matter: an unseen mass component whose
extended halo provides the missing inward acceleration.  This explanation is
internally consistent and broadly accepted, but it depends on a particle species
that has evaded direct detection in three decades of increasingly sensitive
laboratory experiments.  Modified-gravity alternatives, from MOND to $f(R)$
theories, attempt to alter the force law rather than introduce new matter,
but face systematic difficulties with cluster-scale dynamics, the bullet-cluster
mass offset, and the cosmic microwave background power spectrum.

This paper explores a different possibility.  The missing acceleration may arise
not from hidden mass and not from a modified force law, but from
\emph{hidden-dimensional momentum}: a contribution to the radial force budget
that originates in a fifth spatial coordinate $w$ and projects onto the galactic
plane through a logarithmic coupling.  The central translation is:

\begin{keybox}
\[
  \InfoD_i \;\longrightarrow\; \pw^{(i)} = m_i \frac{dw_i}{d\tau} .
\]
The hidden variable is not an inaccessible information debt accumulated over
cosmic history.  It is the current momentum of a particle along a hidden fifth
spatial axis.
\end{keybox}

\noindent
This replacement makes the theory local, kinematic, and directly testable using
existing mass-profile and spectroscopic velocity data.  The paper does not
declare that dark matter does not exist; it constructs a geometric mechanism
capable of producing dark-matter-like rotation signatures, and asks whether
that mechanism can be distinguished from a halo fit using present observations.

\medskip
\noindent
The paper is organised as follows.
Section~\ref{sec:ikk} establishes the framework's identity as Inverse
Kaluza-Klein Scale Symmetry: the compactification rejection, the reversed
inference direction, the five-axiom chain, and the IKK dark-sector
restatement.
Section~\ref{sec:geometry} develops the five-dimensional geometric engine:
the extended coordinate space, five-velocity, five-momentum, the 5D spacetime
interval, and the local frame-dissection equation.
Section~\ref{sec:framework} derives the logarithmic fifth-force potential,
the radial acceleration, the rotation-curve equation, and the geometric lasso.
Section~\ref{sec:fitting} provides the full numerical parameterisation protocol
for regression against data, including the switching function, the transition
radius, ghost-source corrections, and the escalation hierarchy.
Section~\ref{sec:schools} presents three interpretive schools.
Section~\ref{sec:tests} outlines observational tests.
Section~\ref{sec:limits} discusses limitations and falsifiability.
Section~\ref{sec:conc} concludes.

%%═══════════════════════════════════════════════════════════════════════════════
\section{Framework Identity: Inverse Kaluza-Klein Scale Symmetry}
\label{sec:ikk}
%%═══════════════════════════════════════════════════════════════════════════════

\subsection{Naming and Inference Direction}

This framework is called \textbf{Inverse Kaluza-Klein Scale Symmetry} because
it reverses the inference direction of standard Kaluza-Klein theory while
explicitly rejecting the compactification assumption.

In standard Kaluza-Klein theory, an extra spatial dimension is introduced and
compactified --- classically as a small circular coordinate $S^1$ --- and the
four-dimensional force behaviour observed in ordinary spacetime is derived as a
geometric consequence of that compact structure:
\[
  D_{5,\,\text{compact}} \;\xrightarrow{\;\text{KK}\;}\; F_{\text{obs}} .
\]
The geometry comes first; the observed force follows.

The present framework begins at the other end.  It starts with observed
force-scale behaviour, projection failure across measurement scales, and
residual identity content that ordinary observation cannot recover.  From
those residuals it infers what kind of missing degree of freedom must exist:
\[
  F_{\text{obs}},\; R_{\text{obs}}
  \;\xrightarrow{\;\text{IKK}\;}\;
  D_{\text{permission}} ,
\]
where $D_{\text{permission}}$ is \emph{not} a literal curled-up coordinate.
It is a missing degree of freedom required to explain why projection at a
given observational scale loses structure.

\begin{keybox}
\textbf{The key distinction.}\quad
Standard KK: \emph{compact extra geometry $\to$ force behaviour.}\\
Inverse KK: \emph{force-scale residuals $\to$ inferred scale permission.}

The framework is adjacent to Kaluza-Klein in philosophical origin but opposite
in method.  A tiny compact circle is not the theory.  It is the historical
mirror the theory is walking away from.
\end{keybox}

\subsection{Rejection of Compactification}
\label{sec:nocompact}

The coordinate $w$ introduced in this framework should not be interpreted as
\[
  w = \theta R ,\quad \theta \in [0,\,2\pi) ,
\]
nor as $w \in S^1$.  The framework explicitly rejects this identification.
Instead:
\[
  w \;\in\; S_{\text{depth}} ,
\]
where $S_{\text{depth}}$ denotes a scale-depth identity space representing
projection access, unresolved identity content, and hidden residual structure
at a given observational scale.

The distinction is fundamental:

\medskip
\begin{center}
\begin{tabular}{@{}p{0.44\textwidth}p{0.44\textwidth}@{}}
\toprule
\textbf{Compact coordinate} (KK) & \textbf{Scale-depth coordinate} (IKK) \\
\midrule
A spatial loop curled below detection limits &
A bookkeeping axis for unrecovered projection content \\[6pt]
Requires a specific geometric radius $R$ &
Requires only that $P_S(I) \neq I$ \\[6pt]
Fails if the loop has not been found &
Confirmed whenever projection is demonstrably incomplete \\[6pt]
``Something is curled up'' &
``Something is not recovered by the current projection'' \\
\bottomrule
\end{tabular}
\end{center}

\medskip\noindent
The framework therefore does not need the unobserved structure to occupy
microscopic circular geometry.  It only requires that projected observations
fail to recover the full identity content.

\medskip
\begin{keybox}
\textbf{Corrected statement.}\quad
Inverse KK does not compactify an extra dimension.
It infers a missing scale permission from projection residuals.
\end{keybox}

\subsection{Axiom Chain}
\label{sec:axioms}

The following five axioms replace any earlier informal motivation and
constitute the logical skeleton of the framework.

\medskip
\noindent\textbf{Axiom 1: Projection is incomplete.}\\
Observation at scale $S$ does not recover the full identity $I$:
\[
  O_S = P_S(I) , \qquad P_S(I) \neq I ,
\]
unless $P_S$ is perfectly injective and complete, which physical measurement
is generally not.

\medskip
\noindent\textbf{Axiom 2: Residual identity exists.}\\
The residual at scale $S$ is:
\[
  R_S = U_S(I) - E_S(O_S) ,
\]
where $U_S(I)$ is the unprojected identity content available at scale $S$,
and $E_S(O_S)$ is the expected reconstruction from observation alone.
If $R_S \neq 0$, the observed state is not the full state.

\medskip
\noindent\textbf{Axiom 3: Scale permission replaces compactification.}\\
Standard KK uses a compactified geometric extension $D_5 \sim S^1$.
This framework uses scale permission:
\[
  \mathcal{S}_n = (D_n,\; M_n,\; L_n) ,
\]
where $D_n$ is a degree of freedom or permission available at scale $n$,
$M_n$ is the mediator or interaction carrier active at that scale, and
$L_n$ is the observational limit.  The extra structure is not required to
be spatially compact.

\medskip
\noindent\textbf{Axiom 4: Packing creates higher-scale identity.}\\
Components at scale $n$ assemble into identities at scale $n+1$:
\[
  C^{(n)}_{\Pi_n} \;\longrightarrow\; I^{(n+1)} .
\]
This is consistent with the recursive scale ladder: quarks pack into
hadrons, hadrons into nuclei, nuclei into atoms, atoms into molecules,
molecules into materials, materials into gravitational systems, and at each
level new projection residuals appear.

\medskip
\noindent\textbf{Axiom 5: Projection residuals imply hidden-active structure.}\\
A projected channel can cancel while internal structure remains active.
For an identity vector $\mathbf{c}_j$ with charges summing to zero:
\[
  I_j = \boldsymbol{\alpha}^T \mathbf{c}_j = 0 ,
  \qquad
  H_j = \mathbf{c}_j^T W \mathbf{c}_j > 0 .
\]
The hidden-active condition $I_j = 0$, $H_j > 0$ says: apparent neutrality
does not imply internal inactivity.  The neutron is the canonical example:
\[
  +\tfrac{2}{3} - \tfrac{1}{3} - \tfrac{1}{3} = 0 ,
  \quad\text{but}\quad
  \left(\tfrac{2}{3}\right)^2 + \left(\tfrac{1}{3}\right)^2 + \left(\tfrac{1}{3}\right)^2
  = \tfrac{2}{3} \neq 0 .
\]
The same logic applies to galactic dark residuals: the visible baryonic
field appears insufficient, but the full identity space contains active
hidden-axis momentum content that the electromagnetic projection does not
resolve.

\figsuggest{Two-column conceptual diagram: \emph{Left} --- standard KK
inference chain (compact $S^1$ $\to$ gauge field $\to$ observed force);
\emph{Right} --- IKK inference chain (observed residual $\to$ projection
gap $\to$ inferred scale permission $D_{\rm permission}$) --- drawn as
mirror-image arrows to make the inversion visually explicit.}

\subsection{Dark-Sector Restatement}
\label{sec:darksector}

Within the IKK framing, the dark-sector interpretation changes.
The coordinate $w$ is not a compact hidden circle.  It is a scale-depth
coordinate representing identity content that couples differently to
gravitational and electromagnetic projection kernels.

The observed dark residual is:
\[
  \rho_{\text{dark}} = \rho_G - \rho_{\text{EM}} ,
\]
where $\rho_G$ is the mass inferred from gravitational lensing and kinematics,
and $\rho_{\text{EM}}$ is the mass inferred from electromagnetic emission.
In the IKK framework:
\[
  \rho_{\text{dark}}(x,y,z;\,t)
  = \int \rho(x,y,z,w;\,t)
    \bigl[K_G(w) - K_{\text{EM}}(w)\bigr]\,dw ,
\]
where $K_G(w)$ and $K_{\text{EM}}(w)$ are the gravitational and electromagnetic
projection kernels over the scale-depth axis.  The dark regime is defined by
the condition:
\[
  K_G(w) \neq K_{\text{EM}}(w)
\]
for some region of scale-depth $w$.

\begin{keybox}
\textbf{IKK definition of dark matter.}\\
Dark matter is gravitationally visible but electromagnetically unresolved
identity content: a projection gap, not invisible particles.
\end{keybox}

\noindent
This framing is consistent with, and stronger than, the galactic rotation
result derived in the sections that follow.  The fifth-force acceleration
$\afive(r)$ is not an ad-hoc correction; it is the radial projection of
hidden-axis momentum content that the standard electromagnetic projection
$P_{\text{EM}}$ fails to recover.

\figsuggest{Venn diagram: three overlapping regions labelled
``Gravitationally visible'', ``Electromagnetically visible'', and
``Scale-depth resolved''.  The intersection of all three is baryonic
matter.  The region gravitationally visible but not electromagnetically
visible is the dark matter zone --- labelled in this framework as
``$K_G(w) \neq K_{\rm EM}(w)$ region of $w$''.}

%%═══════════════════════════════════════════════════════════════════════════════
\section{Five-Dimensional Geometric Engine}\label{sec:geometry}
%%═══════════════════════════════════════════════════════════════════════════════

\subsection{On Dimensional Order, Scale, and the Necessity of Axes}
\label{sec:dimorder}

\textit{A note before the mathematics.}\enspace
The three spatial axes $x$, $y$, $z$ of ordinary experience are not
intrinsically ordered.  There is no physical reason why $x$ should be
``first'', $y$ ``second'', and $z$ ``third.''  We label them for
bookkeeping.  Any rotation in 3D space mixes them freely, and the laws
of physics are symmetric under that mixing.  The axes are not a
hierarchy; they are a \emph{basis} --- the minimum number of independent
directions required to specify a position or transmit a force in the
space we navigate.

\medskip
\noindent
This raises an open question the moment additional dimensions are proposed:
\emph{what makes a dimension relevant?}  Not its label.  Not its position
in the list.  The question is whether a given axis participates in
\textbf{motion and force transmission} at the scale under consideration.

\medskip
\noindent
The pattern from known physics suggests a scale-dependent answer.  At
molecular scales, three spatial axes plus time are sufficient to account
for all observed forces and displacements.  At nuclear scales, the strong
force introduces \emph{colour charge}: a new internal degree of freedom
that has no geometric axis in ordinary 3-space but behaves like one in the
abstract symmetry space of QCD.  At the electroweak scale, isospin adds
another.  At cosmological scales, dark-energy-like effects suggest that
the effective dimensionality of the force landscape may again change.
In each case the threshold for a new ``effective axis'' is not an ontological
declaration but a statement about \textbf{normalised degrees of freedom that
carry distinguishable momentum or mediate distinguishable forces at that
scale}.

\medskip
\noindent
The present framework makes one specific proposal in this pattern.  At
galactic scales --- where $r \sim 1$--$100\,\mathrm{kpc}$ --- a fifth
spatial axis $w$ becomes dynamically non-negligible.  Below galactic scales
the sigmoid gate $G(r) \approx 0$ effectively removes $w$ from the local
force budget, just as the nuclear degrees of freedom are invisible to a
Newton pendulum experiment.  The fifth axis is not ``more important'' than
$x$, $y$, $z$; it simply requires a \emph{fifth normalised axis} to
describe motion and force at a scale where three axes and one time axis
are no longer sufficient.

\medskip
\noindent
The philosophical analogy with $x$, $y$, $z$: just as there is no preferred
\emph{order} among the three ordinary spatial axes (no law of physics cares
whether you call them 1, 2, 3 or $\alpha$, $\beta$, $\gamma$), there may
be no preferred order among extra spatial dimensions either.  The \emph{number}
of axes active at a given scale is what matters, not their labels or their
position in the list.  This framework requires a fifth axis active at
galactic scales.  Whether it is ``the fifth'' or ``one of many'' is a question
about the global topology of the manifold, not about the local physics.

\figsuggest{Two-panel illustration of scale-dependent dimensional relevance:
\emph{left panel} --- a scale ladder from $10^{-15}$\,m (nuclear) to
$10^{23}$\,m (galactic) with the number of effective dynamical axes annotated
at each scale (3 at human scale, 3$+$colour at nuclear scale, 3$+w$ at
galactic scale); \emph{right panel} --- a ``dimensional activation'' sigmoid
curve showing $G(r)$ as a function of radius, visually tying the mathematical
gate to the conceptual idea of a dimension switching on.}

\figsuggest{Conceptual 3D rendering: start with an $x$--$y$ plane (a flat
galaxy disk); extrude a $z$-axis upward (ordinary height); then draw a
second perpendicular arrow $w$ that is \emph{not} in the plane of the screen
and not aligned with $z$ --- visually suggest the axis is ``into'' a
direction the viewer cannot see, and label it: ``$w$: active only at galactic
scales''.  Emphasise that no one axis is hierarchically superior.}

\subsection{Extended Coordinate Space}\label{sec:coords}

\textit{Non-technical summary.}\enspace
Think of ordinary space as a flat surface.  Every object on it has a position
in the east--west and north--south directions, and a height.  This framework
adds a fourth direction that is not east, west, north, south, up, or down ---
it is a direction you cannot currently point to, but along which particles may
still move.  We call this direction $w$.  The proposal is simply: space has
one more direction than the three we can measure, and that extra direction
matters at galactic scales.

The framework begins from a minimal extension of the standard spacetime
manifold.  A fifth spatial coordinate $w$ is appended alongside the ordinary
three spatial coordinates and physical time:

\[
  \XA = (x^1,\,x^2,\,x^3,\,w,\,T)
       = (\mathbf{r},\,w,\,T) .
\]

The hidden dimension is treated as extended --- not compactified on a
sub-Planck scale --- but currently inaccessible to direct electromagnetic
or nuclear measurement at laboratory energies.  The choice of an
\emph{extended} (rather than Kaluza--Klein compactified) dimension is
deliberate: a large hidden dimension can affect galactic-scale dynamics
without modifying solar-system tests, whereas a Planck-scale compactification
contributes negligibly at all observationally accessible scales.
A particle's full 5D position is:

\[
  \mathbf{X}_i = (\mathbf{r}_i,\,w_i,\,T_i) ,
\]

with ordinary 3D position $\mathbf{r}_i = (x_i,y_i,z_i)$, hidden coordinate
$w_i$, and time coordinate $T_i$.

\figsuggest{Schematic of the extended coordinate space: a 2D ``galaxy disk''
plane floating inside a 3D box, with a fourth arrow labelled $w$ projecting
perpendicular to the box --- illustrating how the hidden axis is orthogonal to
all three ordinary spatial directions.}

\subsection{Five-Dimensional Velocity and Momentum}\label{sec:momentum}

\textit{Non-technical summary.}\enspace
Every moving object carries momentum --- the combination of its mass and
velocity that determines how hard it is to stop.  In ordinary 3D space a
baseball has momentum in the direction it is travelling.  If space has a
hidden fourth direction, a particle can also carry momentum \emph{along that
direction}.  We call this hidden-axis momentum $\pw$.  A particle confined
to ordinary space has $\pw = 0$.  A particle drifting in the hidden direction
has $\pw \neq 0$, and the framework proposes that this hidden momentum acts
as the source of an additional gravitational-like force.  The bigger the
hidden momentum, the stronger the force --- but the coupling grows only
logarithmically, so even very large $\pw$ does not produce runaway forces.

The 5D velocity vector is defined with respect to the local proper-time
parameter $\tau$:

\[
  \UA_i = \frac{d\XA_i}{d\tau}
  = \left(\frac{d\mathbf{r}_i}{d\tau},\;\frac{dw_i}{d\tau},\;\frac{dT_i}{d\tau}\right) ,
\]

and the 5D momentum vector is:

\[
  \PA_i = m_i \UA_i .
\]

In the non-relativistic limit along ordinary spatial directions, the
three-momentum is the standard $\mathbf{p}_i = m_i \mathbf{v}_i$.  The
hidden-axis momentum is:

\begin{keybox}
\[
  \pw^{(i)} \;=\; m_i \frac{dw_i}{d\tau} .
\]
\end{keybox}

\noindent
This is the key new degree of freedom.  It is not a conserved charge added
by hand; it is the component of the existing 5D momentum vector along the
hidden direction.  Any particle that moves in $w$ carries $\pw$; any particle
confined to ordinary 3D space has $\pw = 0$.

\figsuggest{Arrow diagram showing the 5D momentum vector $\PA_i$ decomposed
into an ordinary 3D component $\mathbf{p}_i$ (in the plane of the galaxy disk)
and a perpendicular hidden component $\pw\hat{w}$, analogous to the
decomposition of a velocity vector into horizontal and vertical components.}

\medskip
The full 5D momentum in compact notation:

\[
  \PA_i = (\mathbf{p}_i,\;\pw^{(i)},\;E_i/c) ,
  \quad
  |\PA_i|^2 = |\mathbf{p}_i|^2 + (\pw^{(i)})^2 + \frac{E_i^2}{c^2} .
\]

The hidden kinetic content in the non-relativistic approximation is:

\[
  \Kw = \frac{(\pw)^2}{2m} ,
\]

or relativistically with hidden-axis propagation speed $c_w$:

\[
  E_w = \sqrt{(\pw)^2 c_w^2 + m^2 c_w^4} .
\]

For the first-order force model, the classical $\pw = m\,dw/d\tau$ is
sufficient.

\subsection{The Five-Dimensional Spacetime Interval}\label{sec:interval}

\textit{Non-technical summary.}\enspace
In Einstein's special relativity, a moving clock runs slow --- this is time
dilation, confirmed countless times in particle accelerators and GPS satellites.
The faster you move through ordinary space, the slower you age relative to
someone at rest.  This framework extends that idea: if a particle is also
moving through the hidden $w$-direction, that motion contributes an
\emph{additional} slowing of its local clock.  Two particles with identical
visible speeds could age at different rates if one is also drifting through
the hidden dimension.  This differential aging --- called proper-time
dissection --- is the geometric mechanism underlying the fifth force.

The 5D line element is:

\begin{keybox}
\[
  \dS = -c^2\,dT^2 + d\mathbf{r}^2 + \eta_w\,dw^2 ,
\]
\end{keybox}

\noindent
where $\eta_w$ controls the signature of the hidden dimension.  For an
extended spatial fifth dimension, $\eta_w = +1$ and the metric is:

\[
  \dS = -c^2\,dT^2 + dx^2 + dy^2 + dz^2 + dw^2 .
\]

The local proper-time relation follows from $c^2 d\tau^2 = -\dS$:

\[
  c^2 d\tau^2 = c^2\,dT^2 - d\mathbf{r}^2 - dw^2 .
\]

Dividing through by $dT^2$ and defining the ordinary velocity
$\mathbf{v} = d\mathbf{r}/dT$ and hidden-axis velocity
$v_w = dw/dT$:

\begin{keybox}
\[
  d\tau = dT\sqrt{1 - \frac{|\mathbf{v}|^2}{c^2} - \frac{v_w^2}{c_w^2}} ,
\]
\end{keybox}

\noindent
where $c_w$ is the effective propagation speed scale in the hidden dimension
(not necessarily equal to $c$).  This is the fundamental proper-time relation
for the 5D manifold.  It reduces to the standard Lorentz relation when
$v_w = 0$.  The choice $\eta_w = +1$ (spatial, not time-like hidden axis)
is required for the Lorentz-factor structure in the proper-time step to
remain well-defined; a time-like extra dimension would introduce causality
pathologies outside the scope of this framework.

\figsuggest{Spacetime diagram (Minkowski-style) with an additional $w$-axis:
show two particle worldlines with identical ordinary velocities but different
$v_w$, illustrating that the steeper worldline in $w$ corresponds to a shorter
proper-time accumulation per unit coordinate time $T$.}

\subsection{Frame Dissection: Local Proper-Time Step}\label{sec:frame}

\textit{Non-technical summary.}\enspace
Imagine two identical watches: one handed to a star orbiting normally in a
galaxy's disk, another handed to a star with identical orbital speed but also
drifting slowly in the hidden direction.  After one billion years of coordinate
time, the second watch will have ticked fewer times.  This difference in the
rate of ticking --- called frame dissection here --- is not a metaphor; it
follows directly from the 5D spacetime interval.  The practical consequence is
that the second star experiences a slightly different gravitational bookkeeping
and therefore a different effective inward acceleration.  Summed across billions
of stars in a galaxy's outer disk, this difference produces the observed flat
rotation curve.

For a discrete integration step of coordinate time $\Delta T_0$, the local
proper-time step experienced by particle $i$ is:

\[
  \Delta T_i
  = \Delta T_0 \sqrt{1 - \frac{|\mathbf{v}_i|^2}{c^2} - \frac{v_{w,i}^2}{c_w^2}}
  = \frac{\Delta T_0}{\gamma_{5,i}} ,
\]

where the five-dimensional Lorentz factor is:

\[
  \gamma_{5,i}
  = \left(1 - \frac{|\mathbf{v}_i|^2}{c^2} - \frac{v_{w,i}^2}{c_w^2}\right)^{-1/2} .
\]

The critical consequence is:

\begin{keybox}
\centering
Two particles with identical ordinary velocity $\mathbf{v}_A = \mathbf{v}_B$
but different hidden-axis velocity $v_{w,A} \neq v_{w,B}$ experience
different local frame rates: $\Delta T_A \neq \Delta T_B$.
\end{keybox}

\noindent
This gives $\pw$ a direct geometric role without requiring it to be
a mysterious unmeasurable quantity.  A particle with significant hidden-axis
motion moves more slowly through coordinate time per unit proper time; its
local causal structure is modified relative to an observer confined to
ordinary 3D space.

In low-velocity compact systems --- stellar binaries, planetary orbits,
solar-system tests --- both $|\mathbf{v}| \ll c$ and $v_w \approx 0$, so
$\Delta T_i \approx \Delta T_0$ and the fifth-force contribution is
negligible.  In extended disk systems at galactic radii where orbital
kinematics drive appreciable hidden-axis motion, the frame dissection
diverges from the ordinary Lorentz factor and the fifth-force contribution
becomes significant.

\subsection{The Old Information Variable and Its Geometric Replacement}

\textit{Non-technical summary.}\enspace
An earlier version of this theory used a variable called the ``information
delta'' --- essentially a tally of everything that had ever happened to a
particle since the Big Bang.  The idea was evocative but useless in practice:
you cannot measure the full cosmic history of a star.  The geometric
replacement makes the theory computable.  Instead of asking ``what is this
star's history?'', we ask ``how fast is this star currently moving in the
hidden direction?''.  One question is unanswerable; the other can in principle
be inferred from observable kinematics.  This replacement is not a minor
notational change --- it converts the framework from a philosophical proposal
into a falsifiable physical model.

Earlier formulations of this framework introduced an abstract hidden variable
$\InfoD_i$ (the ``information delta'') representing an inaccessible historical
debt accumulated over a particle's lifetime.  This formulation was useful for
motivation but operationally opaque: it is impossible to compute $\InfoD_i$
for a real galaxy without reconstructing the full thermodynamic history of
every particle since the Big Bang.

The geometric replacement is:

\[
  \InfoD_i \;\equiv\; \mathcal{D}(\pw^{(i)}) ,
\]

where $\mathcal{D}$ maps fifth-axis momentum into apparent information
displacement.  In the simplest instantiation, $\mathcal{D}$ is the
logarithmic coupling:

\[
  \mathcal{D}(\pw) = \ln\!\left(1 + \frac{|\pw|}{p_0}\right) .
\]

The ``information delta'' is not an inaccessible memory file.  It is the
observable shadow of current motion along the hidden fifth spatial direction.
This translation converts a metaphysical historical object into a kinematic
one.  The fifth force depends on what a particle is doing now, not on what
it has done since recombination.


%%═══════════════════════════════════════════════════════════════════════════════
\section{The Logarithmic Fifth-Force Framework}\label{sec:framework}
%%═══════════════════════════════════════════════════════════════════════════════

\subsection{Hidden-Axis Momentum as Fifth-Force Charge}\label{sec:charge}

The central identification of the framework is that $\pw$ acts as a fifth-force
charge.  Rather than a new quantum number or a free scalar field, the charge is
a kinematic quantity that any particle already possesses when it moves in the
hidden direction.  The coupling function is chosen to be logarithmic:

\[
  g(\pw) = \ln\!\left(1 + \frac{|\pw|}{p_0}\right) ,
\]

where $p_0$ is a reference momentum scale.  Three candidate forms are
compared:

\begin{table}[H]
\centering
\caption{Candidate fifth-force coupling functions and their properties.}
\begin{tabular}{@{}llll@{}}
\toprule
Form & $g(\pw)$ & Behaviour at large $|\pw|$ & Behaviour at $|\pw|=0$ \\
\midrule
Linear    & $|\pw|$ & Diverges & $0$ \\
Quadratic & $\pw^2$ & Diverges faster & $0$ \\
Logarithmic & $\ln(1+|\pw|/p_0)$ & Grows slowly & $0$ \\
\bottomrule
\end{tabular}
\end{table}

\noindent
The logarithmic form is chosen for three reasons.  First, it saturates at
large $|\pw|$, preventing high-momentum systems from producing unbounded
fifth-force contributions.  Second, it vanishes continuously at $|\pw|=0$:
particles with no hidden-axis motion contribute no fifth force.  Third, the
form $\ln(1+|\pw|/p_0)$ allows small fifth-axis motions to contribute
meaningfully (the argument exceeds unity only when $|\pw| > p_0$, providing
a natural onset scale) while preventing mathematical divergence for fast-moving
compact objects such as neutron-star binaries.

\subsection{The Logarithmic Potential}\label{sec:potential}

The fifth-dimensional potential is:

\[
  \Phifive(r) = -\Lf \ln\!\left(1+\frac{|\pw|}{p_0}\right)
  \ln\!\left(\frac{r_c}{r+r_c}\right) ,
\]

where $\Lf$ is the fifth-force coupling constant (dimensions of acceleration)
and $r_c$ is a core-softening radius preventing divergence at the origin.
The radial acceleration follows from $\afiver(r) = -d\Phifive/dr$:

\[
  \afiver(r) = \Lf \ln\!\left(1+\frac{|\pw|}{p_0}\right) \frac{1}{r+r_c} .
\]

At large radius ($r \gg r_c$) this behaves as $\afiver \sim 1/r$, precisely
the scaling needed to sustain a flat outer-disk velocity curve.  At small
radius ($r \ll r_c$) the acceleration saturates at $\Lf\ln(1+|\pw|/p_0)/r_c$.

\subsection{Radial Propagation Function and Screening}\label{sec:screening}

A purely global fifth-force profile active at all scales would conflict with
precision gravity tests in the solar system.  The radial propagation function
$\Sfive(r)$ controls the spatial extent of the fifth-force contribution.  The
preferred gated form is:

\[
  \Sfive(r) = \frac{1}{1+e^{-k(r-r_0)}} \cdot \frac{1}{r+r_c} ,
\]

where $r_0$ is the activation radius (of order a few kiloparsecs) and $k$
controls transition sharpness.  The limiting behaviour is:

\[
  r \ll r_0 \;\Rightarrow\; \Sfive \approx 0
  \qquad
  r \gg r_0 \;\Rightarrow\; \Sfive \approx \frac{1}{r+r_c} \approx \frac{1}{r} .
\]

The full fifth-force acceleration vector is then:

\begin{keybox}
\[
  \afive(r) = -\Lf \ln\!\left(1+\frac{|\pw|}{p_0}\right)
  \frac{1}{1+e^{-k(r-r_0)}} \frac{1}{r+r_c}\,\rhat .
\]
\end{keybox}

\noindent
At large disk radii the gate saturates at unity, the core softening becomes
negligible, and $|\afive| \sim \Lf\ln(1+|\pw|/p_0)/r$, giving the $1/r$
profile required for flat rotation curves.  The gate naturally suppresses the
fifth force in compact systems without requiring separate screening mechanisms
for each regime.

\subsection{Gate Benchmark: Molecular Dynamics and Crystal Lattice Consistency}
\label{sec:mdbenchmark}

\begin{keybox}
\textbf{Gate benchmark.}\quad
Crystal-lattice dynamics and classical molecular dynamics (MD) simulations are
completely unaffected by the fifth-force correction---a direct and exact
consequence of $G(r)$ evaluated at atomic and sub-atomic scales.
\end{keybox}

\noindent
The switching function at inter-atomic separations ($r_{\rm atom} \sim
1\,\text{\AA}$--$10\,\text{\AA}$) satisfies
\[
  G(r_{\rm atom})
  = \frac{1}{1+e^{-k(r_{\rm atom}-r_0)}}
  \approx 0 ,
\]
because $k r_0 \gg 1$ and $r_{\rm atom} \ll r_0$.
Consequently the fifth-force acceleration at atomic separations is
\[
  \afiver(r_{\rm atom})
  = -\Lf\ln\!\left(1+\frac{|\pw|}{p_0}\right)
    \underbrace{G(r_{\rm atom})}_{\approx\,0}
    \frac{1}{r_{\rm atom}+r_c}
  \approx 0 .
\]
The same suppression applies at nuclear scales ($r \sim 1\,\text{fm}$),
chemical-bond scales ($r \sim 0.5$--$3\,\text{\AA}$), and planetary scales
($r \lesssim R_\odot$).  The transition radius $r_0$ sits firmly in the
kiloparsec regime; no fifth force leaks into regimes already well-described by
classical or quantum mechanics.

This is an important initial consistency gate for the framework.  Existing
force-field parametrisations (AMBER, CHARMM, OPLS, \emph{ab initio} MD) require
no patching, and QM/MM boundary conditions are unmodified.  Every successful
MD reproduction of crystal phonon spectra, protein folding, and materials
elasticity carries over without modification.  The fifth force is
\emph{galactic in origin, galactic in scope}.

\subsection{The Geometric Lasso}\label{sec:lasso}

The five-dimensional momentum, frame dissection, and fifth-force are linked
in a closed kinematic chain:

\begin{keybox}
\[
  \mathbf{v} \;\longrightarrow\; \pw
  \;\longrightarrow\; \Delta T
  \;\longrightarrow\; \afive
  \;\longrightarrow\; \text{orbital confinement.}
\]
\end{keybox}

\noindent
Explicitly:

\begin{align}
  \pw &= m\,v_w , \\[4pt]
  \gamma_5 &= \left(1 - \frac{v^2}{c^2} - \frac{v_w^2}{c_w^2}\right)^{-1/2} , \\[4pt]
  \Delta T &= \frac{\Delta T_0}{\gamma_5} , \\[4pt]
  \afive &= -\Lf\ln\!\left(1+\frac{|\pw|}{p_0}\right)\frac{1}{r+r_c}\,\rhat .
\end{align}

The circular orbital balance condition $v^2/r = a_N(r) + \afiver(r)$ then
gives the rotation-curve equation.

This chain should not be read as a circular argument --- ``the star stays in
orbit because its orbit generates the force that keeps it in orbit.''
The correct reading is: baryonic orbital motion sources a hidden-axis momentum
footprint; that footprint couples logarithmically to a radial acceleration
field; the radial field adds an inward acceleration term that supplements
the Newtonian term.  The hidden-axis momentum is set by the current kinematic
state, not by a self-referential constraint.

\figsuggest{Flow diagram of the geometric lasso: four nodes ($\mathbf{v}$,
$\pw$, $\Delta T$, $\mathbf{a}_5$) connected by directed arrows, with a
return arrow from $\mathbf{a}_5$ to ``orbital confinement'' labelled
``changes $\mathbf{v}$ next step'' --- making explicit that the chain is
a forward Euler step, not a simultaneous constraint.}

\subsection{Rotation Curve Equation}\label{sec:rotcurve}

For circular motion at radius $r$:

\[
  \frac{v^2(r)}{r} = \frac{GM(r)}{r^2} + \afiver(r) .
\]

Multiplying through by $r$:

\begin{keybox}
\[
  v^2(r) = \underbrace{\frac{GM(r)}{r}}_{\vN^2(r)}
  + \underbrace{\Lf\ln\!\left(1+\frac{|\pw(r)|}{p_0}\right)
    \frac{1}{1+e^{-k(r-r_0)}}\frac{r}{r+r_c}}_{\vfive^2(r)} .
\]
\end{keybox}

\subsection{Galaxy Rotation Curve Result}\label{sec:galresult}

Decomposing explicitly:

\[
  v^2(r) = \vN^2(r) + \vfive^2(r) ,
\]

where the Newtonian component is:

\[
  \vN^2(r) = \frac{GM(r)}{r} ,
\]

and the fifth-force velocity contribution is:

\[
  \vfive^2(r) = \Lf\ln\!\left(1+\frac{|\pw(r)|}{p_0}\right)
  \frac{1}{1+e^{-k(r-r_0)}} \frac{r}{r+r_c} .
\]

In the outer disk, where $r \gg r_c$ and $r \gg r_0$, the gate saturates,
the softening becomes negligible ($r/(r+r_c) \to 1$), and:

\[
  \vfive^2(r) \;\approx\; \Lf\ln\!\left(1+\frac{|\pw(r)|}{p_0}\right)
  \;=\; \text{approximately constant}
\]

provided $|\pw(r)|$ varies slowly across the outer disk.

\begin{keybox}
\textbf{Outer-disk limiting form:}
\[
  v^2(r) \;\approx\; \frac{GM(r)}{r} + \Lf\ln\!\left(1+\frac{|\pw|}{p_0}\right) .
\]
The fifth-force term contributes a nearly flat plateau to the rotation curve.
No dark matter particles are required.
\end{keybox}

\noindent
The fifth-force term $\vfive^2$ behaves as a geometric halo parameter.
It varies from galaxy to galaxy through $|\pw(r)|$, which is set by the
current baryonic kinematic state.  Galaxies with higher hidden-axis momentum
(proxied by their baryonic mass and velocity distribution) will exhibit a
larger fifth-force contribution.  This provides a natural mechanism for
galaxy-to-galaxy diversity without requiring individual dark matter halo fits.

\medskip
Four key behaviours are produced simultaneously by this single equation:

\begin{enumerate}[leftmargin=1.5em]
  \item Near-centre suppression from the sigmoid gate.
  \item Smooth turn-on around $r_0$.
  \item Outer-disk flattening from the $\sim 1/r$ radial profile.
  \item Logarithmic momentum scaling preventing high-$\pw$ runaway.
\end{enumerate}

\figsuggest{Rotation curve plot showing three curves on the same axes: (1)
Newtonian prediction $\vN(r)$ falling after the bulge, (2) the fifth-force
contribution $\vfive(r)$ rising from zero at small $r$ and levelling off at
large $r$, and (3) the combined prediction $v(r) = \sqrt{\vN^2 + \vfive^2}$
approximately flat in the outer disk --- compared to a representative SPARC
data set.}

%%═══════════════════════════════════════════════════════════════════════════════
\section{Numerical Parameterisation and Phenomenological Curve-Fitting}%
\label{sec:fitting}
%%═══════════════════════════════════════════════════════════════════════════════

\noindent
The fifth-force correction is not intended to function as a relabelled dark
matter halo.  In a conventional halo model, the missing acceleration is
attributed to an unseen mass density profile $\rho_{\rm DM}(r)$ whose shape
may be adjusted to match any observed rotation curve.  In the present
framework, the residual acceleration is instead sourced by the hidden-axis
momentum proxy $\pw(r)$, coupled through a logarithmic force law and
constrained radial switching function.  The model must therefore reproduce
galactic residuals from baryonic kinematics and a small set of stable coupling
parameters.  If successful fits require arbitrary galaxy-specific tuning of
ostensibly universal parameters, the framework loses predictive value and
becomes merely a reparameterised halo model.

\subsection{Optimisation Protocol and Parameter Constraints}\label{sec:optim}

The natural benchmark dataset is the SPARC database, which provides baryonic
mass models and high-resolution rotation curves for a diverse sample of nearby
disk galaxies.  We write the total observed radial acceleration as:

\[
  \aobs(r) = \abar(r) + \afiver(r) ,
\]

where $\abar(r) = \vbar^2(r)/r$ is the classical baryonic acceleration and
$\afiver(r)$ is the proposed fifth-force residual.  The empirical residual
target is:

\[
  \Delta a(r) = \aobs(r) - \abar(r)
  = \frac{\vobs^2(r) - \vbar^2(r)}{r} .
\]

The model is fitted by requiring $\afiver(r;\,\Theta) \approx \Delta a(r)$,
where the parameter vector is:

\[
  \Theta = \{\Lf,\,p_0,\,r_c,\,r_0,\,k,\,\alpha_S,\,\beta_S\} .
\]

Here $\Lf$ is the fifth-force coupling constant, $p_0$ is the reference
hidden-momentum scale, $r_c$ is the core-softening radius, $r_0$ is the
transition radius, $k$ controls transition sharpness, and $\alpha_S,\beta_S$
parameterise the extended switching function (Section~\ref{sec:switching}).

The optimisation objective is:

\[
  \chi^2(\Theta) = \sum_{j=1}^{N}
  \left[\frac{\vobs(r_j) - \vmod(r_j;\,\Theta)}{\sigma_{v,j}}\right]^2 ,
\]

with the model velocity:

\[
  \vmod^2(r;\,\Theta)
  = \vbar^2(r) + r\,\afiver(r;\,\Theta)
  = \vbar^2(r) + \Lf\,r\ln\!\left(1+\frac{|\pw(r)|}{p_0}\right)
    \Sfive(r;\,\Theta) .
\]

The decoupling target is:

\begin{keybox}
\[
  \vobs^2(r) - \vbar^2(r) \;=\; r\,\afiver(r) ,
  \qquad \text{i.e.,} \qquad
  \Delta a(r) = \afiver(r) .
\]
This is the clean regression identity.  Every parameter in $\Theta$ is
constrained by the observed minus baryonic velocity squared, not by a
free-form halo profile.
\end{keybox}

\subsection{Non-Monotonic Scaling and the Switching Function}\label{sec:switching}

A monotonic fifth-force profile risks the failure mode common to many
modified-gravity proposals: a correction that works on galactic scales but
leaks into solar-system or cluster regimes where it should be suppressed.
The switching function $S(r)$ modulates the fifth-force field, allowing it
to be attractive, suppressed, or sign-reversed depending on radius.

The full gated acceleration is:

\[
  \afiver(r) = \Lf\ln\!\left(1+\frac{|\pw(r)|}{p_0}\right) \Seff(r)
  \frac{1}{r+r_c} ,
\]

where the sign convention of $\Seff$ determines the regime:

\[
  \Seff(r) > 0 \;\Rightarrow\; \text{attractive / binding regime,} \qquad
  \Seff(r) = 0 \;\Rightarrow\; \text{transition boundary,} \qquad
  \Seff(r) < 0 \;\Rightarrow\; \text{repulsive / expansion-like regime.}
\]

The smooth third-order switching form is:

\[
  S(u) = \tanh\!\left(\alpha_S u + \beta_S u^3\right),
  \qquad
  u = \frac{r - \rcrit}{R_S} ,
\]

producing a continuous transition from $S \to -1$ on one side of $\rcrit$
to $S \to +1$ on the other.  Combined with the sigmoid activation gate:

\begin{keybox}
\[
  \Seff(r) = G(r)\,\tanh\!\left[\alpha_S\!\left(\frac{r-\rcrit}{R_S}\right)
  + \beta_S\!\left(\frac{r-\rcrit}{R_S}\right)^{\!3}\right] ,
  \qquad
  G(r) = \frac{1}{1+e^{-k(r-r_0)}} .
\]
\end{keybox}

\noindent
For galactic-disk binding, the sign convention is $\Seff(r) > 0$ inside the
disk, and the full fifth-force acceleration is:

\[
  \afiver(r) = \Lf\ln\!\left(1+\frac{|\pw(r)|}{p_0}\right)
  \Seff(r) \frac{1}{r+r_c} .
\]

\subsection{Physical Interpretation of the Transition Radius}\label{sec:transition}

The transition radius $\rcrit$ is introduced as a phenomenological placeholder
for a physical phase boundary in the fifth-dimensional slice.  It marks the
radial point where the fifth-force topology changes sign, saturates, or
switches coupling regime.  The preferred terminology throughout the paper is
\emph{transition radius} or \emph{switching radius}; in the vicinity of the
fifth-dimensional phase boundary the coupling changes topology.

Two physical interpretations are supported by the model.

\subsubsection{The Self-Binding Lasso}

Within the galactic disk boundary ($r < \rcrit$), the switching function is
chosen so that $\Seff(r) > 0$.  The fifth force contributes an inward radial
acceleration $\afiver(r) > 0$, and the circular velocity relation becomes:

\[
  \vobs^2(r) = \vbar^2(r) + r\,\afiver(r) .
\]

When $\afiver \sim 1/r$, the term $r\,\afiver$ is approximately constant,
naturally producing a flat outer rotation curve.  In this interpretation,
baryonic motion sources a hidden-axis momentum correction whose radial
projection adds an inward acceleration term.  This is the claim; it is not
the circular argument that ``the star binds itself by its own orbit.''
The hidden-axis momentum is a consequence of the star's current kinematic
state; the force is a consequence of the hidden-axis momentum; the orbit is
a consequence of the combined forces.  The chain is directed, not circular.

\subsubsection{The Inversion Boundary}

At $r = \rcrit$, the switching function crosses zero: $\Seff(\rcrit) = 0$.
Beyond the transition radius, $\Seff(r) < 0$, and the fifth-force contribution
reverses sign:

\[
  \afiver(r) \;\to\; -\afiver(r) .
\]

This creates a transition from inward binding to outward expansion-like
behaviour.  Such a sign flip could be used to model large-scale structure
effects, void boundaries, or cluster-scale residuals.  For the first paper,
the inversion boundary is included as a possible phenomenological feature and
is not declared to be the key to all of cosmology.  The universe has survived
worse marketing.

\subsection{Regression Targets and Falsifiability}\label{sec:falsifiability}

The fitting protocol must test whether the fifth-force parameters are
physically meaningful across galaxy populations, not merely whether residuals
can be minimised for individual galaxies.  Three levels of constraint are
applied.

\textbf{Level A --- Local rotation-curve fit.}
For each galaxy $g$:
\[
  \chi^2_g = \sum_{j=1}^{N_g}
  \left[\frac{\vobs(r_j) - \vmod(r_j)}{\sigma_{v,j}}\right]^2 .
\]

\textbf{Level B --- Population consistency.}
The universal parameters $\{\Lf,\,p_0,\,r_c\}$ must remain approximately
stable across the galaxy sample.  Galaxy-specific quantities
$\{|\pw(r)|,\,r_0,\,\rcrit\}$ may vary, but within bounds constrained by
the observed mass distribution, velocity dispersion, and morphology.  If
$\Lf$ drifts by orders of magnitude between galaxies, the model is
functioning as a disguised halo fit.

\textbf{Level C --- Residual acceleration relation.}
SPARC-based radial-acceleration studies demonstrate a tight empirical relation
between $\aobs$ and $\abar$, with systematic structure below an acceleration
scale of approximately $10^{-10}\,\mathrm{m\,s^{-2}}$.  The fifth-force
model must reproduce this population-level structure, not merely individual
curves.

The model is falsified or materially weakened if $|\pw(r)|$ must be
independently tuned for each galaxy with no relation to observable kinematics,
or if $\Seff(r)$ requires unconstrained galaxy-specific adjustments.  In the
latter case, the framework ceases to be predictive and reduces to a
reparameterised halo model.

\subsection{Resolving Theoretical Leakage: The Ghost-Source Correction}\label{sec:ghost}

The holistic-ledger and temporal-inversion interpretations
(Section~\ref{sec:schools}) are not directly fittable under present observational
constraints.  To preserve predictive utility, their macroscopic effects are
absorbed into a constrained correction term $\Gamma_5(r)$:

\[
  \atot(r) = \asnap(r) + \Gamma_5(r) ,
\]

where $\asnap(r) = \Lf\ln(1+|\pw(r)|/p_0)\,\Seff(r)/(r+r_c)$ is the
kinematic snapshot term.  The correction $\Gamma_5(r)$ is not allowed to
become a free dark-matter halo; it is restricted to one of two paradigms.

\subsubsection{Paradigm A: Co-Local, Co-Sourced Emission}

In the co-local paradigm, the correction term is anchored to the same
baryonic source and hidden-momentum footprint that drives the primary fifth
force:

\[
  \Gamma_5^A(r) = \epsilon_A
  \left[\ln\!\left(1+\frac{|\pw(r)|}{p_0}\right)\right]^{n_A}
  \frac{1}{(r+r_A)^{m_A}} ,
\]

where $\epsilon_A$ is a small leakage amplitude, $n_A$ controls non-linear
momentum response, $r_A$ is a softening scale, and $m_A$ controls radial
decay.  This version assumes that any ghost-like residual remains local,
co-sourced, and tied to visible moving matter.  Its advantage is computational
cleanliness; its weakness is difficulty with systems where the gravitational
residual is spatially offset from the baryonic mass.

\subsubsection{Paradigm B: Non-Local Background Decoupling}

In the decoupled paradigm, the correction term represents an ambient
fifth-dimensional background configuration:

\[
  \Gamma_5^B(r) = \epsilon_B\,\Bfield(r,T)\,\frac{1}{r+r_B} ,
\]

where $\Bfield(r,T)$ is a background metric-intensity field parameterised as:

\[
  \Bfield(r,T) = B_0 + B_1\ln\!\left(1+\frac{r}{R_B}\right) + B_T T(t) .
\]

This allows the fifth-force residual to decouple from local mass distributions.
Its advantage is flexibility for asymmetric or displaced residual structures
such as cluster-scale offsets.  Its disadvantage is a dramatically expanded
parameter space.  Paradigm B is therefore activated only after Paradigm A
fails under well-defined residual tests.

\subsection{Recommended Numerical Hierarchy}\label{sec:hierarchy}

To keep the theory testable, a strict escalation ladder governs which model
complexity is deployed.  Higher levels are not activated until lower levels
have been fitted and found insufficient:

\begin{table}[H]
\centering
\caption{Model escalation hierarchy.}
\label{tab:hierarchy}
\begin{tabular}{@{}cll@{}}
\toprule
Level & Model & Fifth-force term \\
\midrule
1 & Baseline snapshot & $\Lf\ln(1+|\pw|/p_0)/(r+r_c)$ \\
2 & Screened snapshot & $\times\,G(r)$ \\
3 & Switching snapshot & $\times\,\Seff(r)$ \\
4 & Co-local ghost correction & $+\,\Gamma_5^A(r)$ \\
5 & Non-local background correction & $+\,\Gamma_5^B(r)$ \\
\bottomrule
\end{tabular}
\end{table}

\noindent
The fitting hierarchy is designed to prevent the fifth-force model from
becoming an unconstrained residual engine.  The kinematic snapshot model
is tested first, using only the logarithmic hidden-momentum force and its
constrained radial switching behaviour.  Historical and geometric leakage
terms are activated only if coherent residual structures remain after the
baseline model has been fitted.  In this way the framework remains falsifiable:
it must either reproduce observed galactic acceleration profiles with stable
parameters, or fail --- without retreating into arbitrary correction fields.

\figsuggest{Radial profile plot of $\Seff(r)$ vs.\ $r$ for a representative
galaxy, showing the sigmoid activation gate $G(r)$ rising from 0 to 1 around
$r_0$, the tanh switching function crossing zero at $r_{\rm crit}$, and the
combined $\Seff(r)$ profile --- illustrating how three distinct radial regimes
(suppressed / binding / reversing) emerge from two smooth functions.}

\figsuggest{SPARC residual scatter plot: observed minus baryonic acceleration
$\Delta a(r)$ vs.\ baryonic acceleration $a_{\rm bar}$ for a sample of SPARC
galaxies, overlaid with the Level-2 screened-snapshot fifth-force prediction
$\afiver(r)$ --- a direct analogue of the McGaugh et al.\ (2016)
radial-acceleration-relation plot.}

\medskip
\noindent\textbf{A note on dimension stacking and temporal independence.
If the hidden axis $w$ is stacked over the time axis $T$ --- that is, if
$w$ and $T$ are treated as parallel dimensions in a higher-dimensional
block rather than as a sequence --- then the $w$-coordinate and its
associated scalar force coupling $g(\pw)$ are likely to be completely
independent of $T(t)$.  A force or momentum contribution that is genuinely
orthogonal to the time dimension carries no natural parametric dependence
on the current value of $T$.  As a direct consequence, $|\pw(r)|$ evaluated
across a galaxy sample will likely not curve-fit correctly to any known
temporal or evolutionary function: it will appear, to observers embedded in
the $T$-dimension and sampling the universe at a single temporal slice, as
genuinely random variation.  The scatter in fifth-force proxy values between
galaxies at the same mass and morphological class may therefore not be noise
to be removed, but signal to be respected --- a direct observational signature
of a force coordinate that does not know what time it is.}

\medskip

%%═══════════════════════════════════════════════════════════════════════════════
\section{Three Schools of Interpretation}\label{sec:schools}
%%═══════════════════════════════════════════════════════════════════════════════

The hidden variable $\pw$ admits three distinct physical interpretations.
These are presented as competing schools, not simultaneous claims.  The
paper adopts School~I as the working framework; Schools~II and~III are
retained as theoretical alternatives that make contact with deeper geometric
or thermodynamic structures.

\subsection{School I: The Kinematic Snapshot}\label{sec:school1}

\textbf{Core axiom.} The fifth force depends strictly on the particle's
\emph{current} fifth-axis momentum, measured within a finite local frame
dissection.  The force carries no historical memory:

\[
  \afive(t_n) = \afive\!\left(\pw(t_n),\,\mathbf{r}(t_n)\right) ,
  \quad \text{not} \quad
  \afive\!\left(\int_0^{t_n} \pw(t)\,dt\right) .
\]

\textbf{Field mechanic.} At each discrete frame step $\Delta T_n$, the
particle emits a fifth-force field determined by its instantaneous hidden-axis
momentum:

\[
  q_5(t_n) = \ln\!\left(1+\frac{|\pw(t_n)|}{p_0}\right) ,
\]

\[
  \mathbf{a}_5^{(n)} = -\Lf\,q_5(t_n)\,G(r^{(n)})\,\frac{1}{r^{(n)}+r_c}\,\rhat^{(n)} .
\]

\textbf{Practical advantage.} This school is local, deterministic, and
computationally tractable.  The hidden-axis momentum $\pw(r)$ is treated as
a proxy derived from baryonic kinematics:

\[
  v^2(r) = \frac{GM(r)}{r} + \Lf\ln\!\left(1+\frac{|\pw(r)|}{p_0}\right)
  \frac{1}{1+e^{-k(r-r_0)}} \frac{r}{r+r_c} .
\]

Astronomers can fit the model using current observational data: $M(r)$ from
luminosity and gas mapping, $v(r)$ from Doppler spectrometry, and a fitted
fifth-momentum proxy $|\pw(r)|$ constrained to correlate with baryonic
momentum structure.

\textbf{Observational prediction.} The fifth-force contribution should
correlate with current kinematic measurements --- $v(r)$, $M(r)$,
$\sigma_v(r)$ --- rather than formation history, metallicity, or AGN activity.
If a star is perturbed by a passing body, the fifth-force contribution changes
in the next frame step because $\pw(t_n) \to \pw(t_{n+1})$.  The model is
therefore highly responsive to current dynamics.

\textit{This is the strongest working school.}  It is adopted as the primary
framework throughout this paper.

\subsection{School II: The Holistic Ledger}\label{sec:school2}

\textbf{Core axiom.} The fifth force depends on a particle's cumulative
lifetime energy, entropy, and information delta:

\[
  Q_{5,i} = \int_{\Gamma_i} dD_i
  = \int_0^T \!\!\left(\alpha_E\,dE_i + \alpha_S\,dS_i + \alpha_I\,d\InfoD_i\right) ,
\]

where $Q_{5,i}$ is the lifetime fifth-dimensional charge.

\textbf{Field mechanic.} Every collision, emission, absorption, decay,
binding transition, and annihilation modifies the particle's historical
ledger:
\[
  \afive = -\Lf\,f(Q_5)\,\Sfive(r)\,\rhat ,
  \quad
  f(Q_5) = \ln\!\left(1+\frac{|Q_5|}{Q_0}\right) .
\]

\textbf{Observational prediction.} Two galaxies with identical visible mass
and current velocity distributions may exhibit different rotation profiles
if their thermodynamic histories differ.  A galaxy with a violent merger or
starburst history ($Q_{5,A} \gg Q_{5,B}$) could show a stronger fifth-force
contribution than a quiescent galaxy, even if their present-day baryonic
distributions are similar.  This school predicts correlations with star
formation history, metallicity, AGN activity, and stellar age dispersion.

\textbf{Strength.} Explains galaxy-to-galaxy rotation diversity through
formation history rather than requiring galaxy-specific halo tuning.

\textbf{Weakness.} Computing $Q_{5,i}$ directly requires reconstructing the
full thermodynamic history of every relevant particle since recombination.
Scientifically interesting; operationally very difficult to constrain.

\subsection{School III: The Temporal Inversion}\label{sec:school3}

\textbf{Core axiom.} The fifth dimension $w$ is more fundamental than time $T$.
Time is not a primary driver; it is a lower-dimensional projection of a fixed
higher-dimensional geometry:

\[
  T = \Pi_T(w,\,x,\,y,\,z) ,
\]

where $\Pi_T$ is a projection from the 5D architecture into experienced 4D
spacetime.  The hierarchy is:

\[
  \text{5D geometry} \;\longrightarrow\; \text{4D spacetime projection}
  \;\longrightarrow\; \text{observed temporal evolution.}
\]

\textbf{Field mechanic.} Particles do not evolve through time in the ordinary
sense.  Their apparent time evolution is caused by the 4D universe slicing
through a fixed 5D structure.  The fifth force is a geometric track, not a
dynamical emission:

\[
  \afive = -\nabla_r \Phifive(w,\,r) ,
  \quad
  \Phifive(w,r) = -\Lf(w)\ln\!\left(\frac{r_c}{r+r_c}\right) .
\]

The observed fifth force $\afiver = \Lf(w)/(r+r_c)$ is a consequence of the
fixed 5D geometry, not of any dynamical emission process.

\textbf{Observational prediction.} Rotation-curve flattening, cosmic
expansion anomalies, and large-scale structure residuals all become
cross-section effects: observed dynamics are projections of a static 5D
geometry.  Galaxies rotate along pre-existing fifth-dimensional tracks.

\textbf{Strength.} Most philosophically powerful and geometrically radical.
Connects the dark-energy problem, the rotation-curve problem, and the
large-scale structure problem through a single 5D geometric origin.

\textbf{Weakness.} Least accessible experimentally.  Requires a clear mapping
$w(r) \to v(r)$ or $\Phifive(w,r) \to$ observable residuals to be falsifiable.
Without such a mapping it becomes cathedral architecture with equations.

\subsection{Comparative Summary}\label{sec:synthesis}

\begin{table}[H]
\centering
\caption{Comparison of the three interpretive schools.}
\label{tab:schools}
\begin{tabularx}{\textwidth}{@{}lXXX@{}}
\toprule
Metric & School I: Kinematic Snapshot & School II: Holistic Ledger & School III: Temporal Inversion \\
\midrule
Primary variable & Current $\pw(t_n)$ & Lifetime integral $Q_{5,i}$ & Fixed 5D coordinate $w$ \\[4pt]
Force source & Instantaneous momentum footprint & Accumulated history & Static 5D geometry \\[4pt]
Time role & Discrete frame parameter $\Delta T$ & Chronological aggregator & Projection beneath $w$ \\[4pt]
Rotation-curve fit & Direct and tractable & Difficult to constrain & Requires 5D map inversion \\[4pt]
Galaxy diversity mechanism & $|\pw(r)|$ varies with kinematics & Formation-history variation & Geometric track variation \\[4pt]
Main strength & Testable with current data & Explains scatter in DM-analogue & Deepest geometric interpretation \\[4pt]
Main weakness & Needs $\pw$ proxy from data & History intractable to reconstruct & Risks unfalsifiability \\[4pt]
Recommended role & Primary framework & Speculative extension & Theoretical boundary case \\
\bottomrule
\end{tabularx}
\end{table}

\noindent
The present work adopts the kinematic snapshot interpretation because it is
local, calculable, and directly testable using current galactic velocity data.
Schools~II and~III are retained as theoretical alternatives but are not
required for the first-order rotation curve model.

\figsuggest{Radar / spider chart comparing the three schools across six axes:
(1) testability with current data, (2) parameter efficiency, (3) ability to
explain galaxy-to-galaxy scatter, (4) connection to deeper geometry,
(5) computational tractability, (6) falsifiability score --- visually
summarising the trade-offs in Table~\ref{tab:schools}.}

%%═══════════════════════════════════════════════════════════════════════════════
\section{Observational Tests}\label{sec:tests}
%%═══════════════════════════════════════════════════════════════════════════════

\subsection{Rotation Curve Fitting Against SPARC}\label{sec:sparc}

The primary observational test is direct regression of the model velocity:

\[
  \vmod^2(r) = \vbar^2(r) + \Lf\ln\!\left(1+\frac{|\pw(r)|}{p_0}\right)
  \frac{1}{1+e^{-k(r-r_0)}} \frac{r}{r+r_c}
\]

against the SPARC rotation-curve database.  The baryonic contribution
$\vbar^2(r)$ is taken from the provided stellar mass models and gas
distributions, leaving $\Theta = \{\Lf,\,p_0,\,r_c,\,r_0,\,k\}$ plus a
proxy for $|\pw(r)|$ as the fitting parameters.

The hidden-axis momentum proxy is the key modelling choice.  Three
physically motivated proxy forms are:

\begin{enumerate}[leftmargin=1.5em]
  \item \textbf{Circular velocity proxy:} $|\pw(r)| \propto m_* v_{\rm circ}(r)$,
    where $m_*$ is a reference stellar mass.  Simple; ties $\pw$ to the
    observed velocity field directly.
  \item \textbf{Surface density proxy:} $|\pw(r)| \propto \Sigma(r) v_{\rm circ}(r)$,
    where $\Sigma(r)$ is the baryonic surface density.  Weights by the local
    mass concentration.
  \item \textbf{Momentum flux proxy:} $|\pw(r)| \propto \rho(r) v(r)$,
    the local three-momentum density.  Connects to the fluid-momentum
    interpretation of hidden-axis sourcing.
\end{enumerate}

Consistency of $\{\Lf,\,p_0,\,r_c\}$ across the sample is the primary
population test.  These parameters should not drift by more than a factor
of a few between galaxies spanning an order of magnitude in mass.

\figsuggest{Scatter plot of fitted $\Lf$ values vs.\ galaxy stellar mass
$M_*$ across a sample of 50 SPARC galaxies, with error bars from bootstrap
fitting --- testing whether $\Lf$ is consistent with a single universal value
or shows systematic mass dependence.}

\subsection{Radial Acceleration Relation}\label{sec:rar}

The empirical radial acceleration relation (McGaugh et al.\ 2016; Lelli et
al.\ 2017) establishes a tight correlation between $\aobs$ and $\abar$
across the SPARC sample.  The fifth-force model predicts:

\[
  \aobs = \abar + \afiver(\abar,\,r) ,
\]

where $\afiver$ is sourced by $|\pw(r)|$, itself a function of the baryonic
momentum distribution.  The model should therefore reproduce the tight
$\aobs$--$\abar$ relation without explicit tuning, provided $|\pw(r)|$
correlates sufficiently tightly with $\abar$.  If the predicted scatter in
the relation exceeds the observed scatter, the proxy model for $|\pw(r)|$
requires refinement.

\subsection{Screening Tests in Compact Systems}\label{sec:screentest}

The sigmoid gate $G(r)$ suppresses the fifth force for $r < r_0$.  Precision
solar-system tests (lunar laser ranging, planetary ephemerides) constrain any
fifth-force contribution at AU scales to below $10^{-13}\,\mathrm{m\,s^{-2}}$.
This requires $G(r_{\rm AU}) \lesssim 10^{-3}$, which is satisfied provided
$r_0 \gg 1\,\mathrm{AU}$ and $k$ is large enough.

Binary-pulsar timing sets complementary constraints at $\sim 1\,\mathrm{AU}$
scale.  The fifth-force contribution to the orbital decay must be negligible
compared to the gravitational-wave prediction from general relativity.

Dwarf spheroidal galaxies provide a middle-scale test: they span radii of
$\sim 1$--$10\,\mathrm{kpc}$ and have very low surface brightnesses.  If
the fifth-force contribution scales with baryonic momentum density, dwarfs
should show less pronounced rotation-curve flattening than high-surface-
brightness spirals.  This is consistent with observations and serves as
a non-trivial prediction of the framework.

\subsection{Population Tests and Morphological Correlations}\label{sec:poptest}

School~I predicts that the fifth-force contribution correlates with
\emph{current} kinematic properties: surface brightness, velocity dispersion,
and baryonic momentum flux.  School~II predicts that it correlates with
\emph{historical} properties: metallicity, star formation rate history, and
merger history.  These predictions are distinct and can in principle be
distinguished using multi-wavelength photometry and resolved stellar population
data.

A practical test: in a sample of galaxies matched in total mass and present-day
velocity dispersion, Schools~I and~II predict the same fifth-force contribution
(since current kinematics are matched) unless merger history is significantly
different.  If residuals from the rotation-curve fit correlate with merger
indicators (asymmetry, colour gradients, disturbed morphology), School~II
gains evidence.  If they do not, School~I is preferred.

%%═══════════════════════════════════════════════════════════════════════════════
\section{Limitations and Falsifiability}\label{sec:limits}
%%═══════════════════════════════════════════════════════════════════════════════

\subsection{The Hidden-Momentum Proxy Problem}

The $\pw(r)$ proxy is the central modelling uncertainty.  The framework
specifies that $\pw$ is the particle's current hidden-axis momentum, but it
does not derive $\pw$ from first principles.  In practice, $|\pw(r)|$ is
treated as a free radial function constrained by the fitting protocol.  If
a successful fit requires a $|\pw(r)|$ profile with no systematic relation to
any observable --- mass, velocity, surface density --- then the model is using
$\pw$ as a hidden halo in all but name.

Falsification criterion: \emph{the model is disqualified if fits to a sample
of $N_{\rm gal} \geq 20$ galaxies require per-galaxy $|\pw(r)|$ profiles with
no mutual correlation across the sample, after controlling for total mass and
morphology.}

\subsection{Parameter Stability Under Sample Variation}

The universal parameters $\{\Lf,\,p_0,\,r_c\}$ must remain stable.  If
they drift significantly between galaxy morphological types (spirals vs.\
dwarfs vs.\ ellipticals), the universality claim fails and the framework
reduces to a flexible fitting function rather than a physical model.

Falsification criterion: \emph{the model is weakened if $\Lf$ varies by
more than a factor of 10 across a morphologically diverse sample without
a physical explanation tied to a fifth-dimension coupling mechanism.}

\subsection{Cluster-Scale Consistency}

Galaxy clusters provide an additional constraint.  The bullet cluster
demonstrates a spatial offset between the baryonic mass (X-ray gas) and the
gravitational lensing mass.  Any geometric fifth-force framework must explain
why the lensing signal follows the collisionless galaxies, not the gas, even
though the gas carries more baryonic mass and --- under most proxy models ---
more hidden-axis momentum.

The framework can accommodate this if $|\pw(r)|$ is sourced primarily by
stellar (collisionless) matter rather than gas, but this requires a specific
physical mechanism --- yet to be derived --- distinguishing the two baryonic
components in their coupling to the fifth axis.

\subsection{Comparison with Halo Models}

The framework is weakened but not falsified if it cannot outperform an NFW
dark matter halo in reduced-$\chi^2$ fits with comparable parameter counts.
It is falsified if it cannot match halo-model fits even with more parameters.
The goal of the first paper is not to defeat dark matter; it is to demonstrate
that the logarithmic fifth-force framework is not immediately excluded by
existing rotation-curve data.

%%═══════════════════════════════════════════════════════════════════════════════
\section{Conclusion}\label{sec:conc}
%%═══════════════════════════════════════════════════════════════════════════════

This paper has developed a fifth-dimensional momentum framework for galactic
rotation curves.  The central replacement is:

\begin{keybox}
\[
  \InfoD_i \;\longrightarrow\; \pw^{(i)} = m_i\frac{dw_i}{d\tau} .
\]
The fifth force is not driven by historical information.  It is driven by
current hidden-axis momentum.
\end{keybox}

\noindent
The key results of the framework are as follows.

The five-dimensional spacetime interval $dS_5^2 = -c^2\,dT^2 + d\mathbf{r}^2 + dw^2$
implies a proper-time dissection $d\tau = dT\sqrt{1 - v^2/c^2 - v_w^2/c_w^2}$
that gives $\pw$ a direct geometric role.  Particles with identical visible
kinematics but different hidden-axis motion experience different local frame
rates.

The fifth-force acceleration follows a logarithmic coupling law:
\[
  \afive(r) = -\Lf\ln\!\left(1+\frac{|\pw|}{p_0}\right)
  \frac{1}{1+e^{-k(r-r_0)}} \frac{1}{r+r_c}\,\rhat ,
\]
which produces a $1/r$ radial profile at galactic scales, $1/r$ being
precisely the scaling needed to sustain flat rotation curves.

The full rotation curve equation is:
\[
  v^2(r) = \frac{GM(r)}{r}
  + \Lf\ln\!\left(1+\frac{|\pw(r)|}{p_0}\right)
    \frac{1}{1+e^{-k(r-r_0)}} \frac{r}{r+r_c} ,
\]
which naturally suppresses the fifth force in compact low-velocity systems
and produces an approximately flat contribution to $v^2(r)$ at outer disk
radii.

A numerical parameterisation protocol has been developed for regression against
the SPARC rotation-curve database, with an escalation hierarchy from baseline
snapshot through co-local and non-local ghost-source corrections.  Population
consistency tests --- requiring that universal parameters $\{\Lf,\,p_0,\,r_c\}$
remain stable across galaxy samples --- provide a clear falsifiability criterion
that distinguishes the framework from a disguised halo fit.

Three interpretive schools have been identified.  The kinematic snapshot school
is adopted as the primary working model because it is local, calculable, and
directly testable using current mass-profile and spectroscopic velocity data.
The holistic ledger and temporal inversion schools are retained as theoretical
extensions for future development.

The framework does not claim to have solved the dark matter problem.  It
proposes a competing geometric mechanism capable of producing dark-matter-like
rotation signatures from hidden-dimensional momentum, written in a form
amenable to direct observational test.

\bigskip

The open question of dimensional order and the independence of the $w$-axis
from the time coordinate $T$ remain theoretical frontiers.  If the $w$-axis
is genuinely orthogonal to $T$ in a block-universe topology, then the scatter
of fifth-force proxy values across galaxy populations may not be noise but
signal: a direct observational signature of a force coordinate that does not
experience time.  Some residuals of the fit may therefore be fundamentally
irreducible --- not a failure of the model but a statement about the geometry
of the universe.

Within the Inverse Kaluza-Klein framing, the framework's position is now
unambiguous: it is not a compactification theory.  It is a projection-residual
inference.  The $w$-coordinate is not a small circle hiding below experimental
reach; it is the bookkeeping axis for identity content that ordinary
four-dimensional projection fails to recover.  The rotation-curve residual is
one instance of that failure.  It will not be the last.

\medskip\noindent
The rotation curve does not lie.  Something is missing from the standard
accounting.  This framework offers one geometric answer.

\end{document}
